% SPDX-License-Identifier: GPL-3.0-or-later
% Copyright (C) 2026 Alessandro Greco
\chapter{Tor 学术研究}\label{cap:29-ricerca}
\aggiornatoal{2026 年 9 月 23 日}{本章引用的研究论文；Tor Project 研究门户；2026 年 9 月 23 日的共识}

\begin{obiettivi}
学完本章后，你将能够：
\begin{itemize}[nosep]
  \item 找到 Tor 研究的重要发表渠道，以及 Tor Project 提供的资源；
  \item 从熵到首次失守时间，定量衡量匿名性；
  \item 选择合适的实验方法：分析建模、模拟、仿真或实际测量；
  \item 使用真实共识模拟数千个客户端在数月内的活动（实验~19）。
\end{itemize}
\end{obiettivi}

从一开始，Tor 的发展就与学术研究紧密相连。洋葱路由诞生于研究实验室（\cref{cap:03-nascita-tor}）；Tor 的设计曾在学术会议发表~\cite{dingledine2004tor}；本书介绍的众多机制，从入口守卫到拥塞控制，都是在回应已发表的研究成果。本章面向希望自己开展研究的读者：到哪里发表成果、如何衡量匿名性、应该选择哪些实验工具，以及必须遵守哪些安全原则。

% ---------------------------------------------------------------------------
\section{Tor 研究发表在哪里？}\label{sec:ric-dove}

匿名网络研究主要发表在两类渠道：
\begin{description}
  \item[Privacy Enhancing Technologies Symposium（PETS）] 以及对应的期刊 \emph{Proceedings on Privacy Enhancing Technologies}（PoPETs）。这是隐私技术领域的专门发表平台；下文讨论的匿名性测量奠基研究发表于 2002 年的 PET 研讨会，早于 PoPETs 期刊创办~\cite{diaz2002measuring}。
  \item[重要安全会议] 包括 USENIX Security、ACM CCS、IEEE Symposium on Security and Privacy，以及 NDSS。\cref{cap:26-attacchi} 引用的多篇论文发表于这些会议。
\end{description}
Free Haven 的 \emph{Anonymity Bibliography} 已经收录这一领域二十多年的研究成果，是文献综述的自然起点。

Tor Project 还维护研究门户，提供数据、工具、技术报告、研究选题以及 Research Safety Board 的相关信息。其建议强调，协议和网络一直在变化：如果不了解内部工作方式就进行测量，很容易得出错误结论。因此，研究人员应尽早与项目联系，了解同时进行的实验和计划中的修改。

% ---------------------------------------------------------------------------
\section{如何衡量匿名性？}\label{sec:ric-metriche}

在 \cref{cap:01-anonimato} 中，我们用某个集合内的不可区分性来定义匿名性。然而，从攻击者的视角看，集合中每个成员成为目标的可能性并不相同。观察流量后，攻击者可能认为某些用户的概率远高于其他人。仅仅知道集合中有多少人并不够。

\paragraph{匿名程度。}2002 年，Díaz 等人提出通过攻击者赋予潜在发送者的概率分布的熵来衡量匿名性~\cite{diaz2002measuring}。如果 $p_i$ 是攻击者认为用户 $i$ 为发送者的概率，熵为
\[ H(X) = -\sum_{i=1}^{N} p_i \log_2 p_i, \]
当所有 $N$ 个用户的概率相同时，熵达到最大值 $H_M = \log_2 N$。\textbf{匿名程度}定义为
\[ d = \frac{H(X)}{H_M}, \qquad 0 \le d \le 1: \]
当 $N\ge 2$ 且后验分布均匀时，$d=1$（这不一定表示攻击者没有获得新信息）；只有一名发送者的概率为 1 时，$d=0$。若 $N=1$，由于 $H_M=0$，此比值无定义。

理解熵的一种直观方式是计算 $2^{H(X)}$：它表示在相同不确定性下，概率相等的用户数量。假设匿名集有 1\,000 人，但攻击者已将嫌疑范围缩小到两人，各有 0.5 的概率，那么不确定性就与只有两人的集合相同：$H=1$ 比特。

\paragraph{单次测量的局限。}熵描述单次观察，但正如 \cref{sec:att-correlazione} 所示，Tor 的暴露风险会随时间累积：每条新电路都为攻击者带来一次新机会。因此 Johnson 等人提出\emph{纵向}指标，包括在特定时间内至少一次失去匿名性的概率，以及首次失守时间的中位数。这些指标可通过结合真实网络数据模拟用户行为计算~\cite{johnson2013users}。下面的实验采用类似方法。

% ---------------------------------------------------------------------------
\section{四种研究方法}\label{sec:ric-metodi}

\begin{description}
  \item[分析建模] 建立数学模型，例如匿名程度模型，或 \cref{sec:att-sybil} 中的 Sybil 攻击模型。它速度快、适用范围广，但依赖各种假设，而实际系统总会在一定程度上偏离假设。
  \item[行为模拟] 使用真实数据模拟系统某个简化的部分，例如路径选择。Johnson 等人开发的路径模拟器 TorPS~\cite{johnson2013users} 和本书实验~19 都采用这种方式。其运行速度快，却无法涵盖模型遗漏的因素。
  \item[网络仿真与模拟] 在本书贯穿始终使用的 Chutney 私有网络中运行真实软件，或者在 Shadow 等离散事件模拟器内运行它~\cite{jansen2012shadow}。Shadow 可以通过虚拟时间，在一台机器上执行数千台中继和客户端（\cref{sec:rp-shadow}）。这种方法尤其适合性能研究，因为新的协议设计可以先在逼真的网络中测试，再部署到公共网络。
  \item[公共网络测量] 只有直接测量才能反映真实网络正在发生的情况，但它对真实用户也有最高风险，因此需要严格的实验设计（\cref{sec:ric-sicurezza}）。
\end{description}
优秀研究通常结合多种方法：数学分析产生假设，大规模模拟检验假设，精心设计的测量再考察模型是否符合现实。

\begin{approfondimento}[在不观察个体的情况下测量]
Jansen 和 Johnson 开发的 PrivCount 展示了测量与隐私保护可以共存~\cite{jansen2016privcount}。参与测量的中继虽然会统计相关事件，但计数器使用与其他参与者共享的随机值进行掩蔽，并依据\emph{差分隐私}原则加入噪声。最终可见的只是经过聚合且刻意带有噪声的结果，从而限制人们推断单个用户信息的能力。2016 年，研究人员利用 PrivCount 估计约有 710\,000 名同时在线客户端，其中约 550\,000 名在某一时刻处于活跃状态；流向 Web 端口的流量占参与测量的出口所记录数据字节的 91\%。这些是 2016 年研究的历史数据，并非当前网络统计。这比 \cref{sec:metrics-utenti} 中通过舍入中继统计保护隐私的做法更进一步。
\end{approfondimento}

% ---------------------------------------------------------------------------
\section{实验：首次失守需要多长时间？}\label{sec:lab-ricerca}

\begin{laboratorio}[19]{模拟数月的电路}
\textbf{目标。}使用真实共识估算：攻击者控制一定比例带宽时，分别在使用与不使用持久守卫的情况下，随时间能使多少客户端失去匿名性。\\
\textbf{前置条件。}实验~7 和实验~15；无须运行私有网络。\\
\textbf{预计时间。}1--2 小时。
\end{laboratorio}

\file{labs/zh-Hans/代码/torlab/simulation.py} 模块复用实验~7 的带宽权重。攻击者增加一台守卫和一台出口，两者各占对应位置加权带宽的比例 $q$。我们模拟 2\,000 名客户端运行 180 天，每个客户端每天建立 20 条电路。

本模型按照\emph{已确认}（confirmed）守卫的近似寿命安排持久守卫的轮换，而不是完整复现实际客户端行为。根据 guard-spec 和 \file{entrynodes.c} 中的代码，已确认守卫在确认后达到 \texttt{GUARD\_CONFIRMED\_MIN\_LIFETIME}（60 天）时会从样本中移除，而 \texttt{CONFIRMED\_ON\_DATE} 会随机向前回溯最长 \texttt{GUARD\_LIFETIME}/10（12 天）。因此，模拟中的轮换间隔为 48--60 天。120 天的 \texttt{GUARD\_LIFETIME} 用于管理未确认守卫在候选样本中的到期时间。真实客户端还会考虑守卫样本和确认列表、可达性及替换策略，无法直接用这一间隔预测。实验的早期版本误将守卫轮换间隔设为 120 天，从而在这个简化模型中低估了结果。作为对照，我们还模拟每条电路都重新选择守卫的客户端。当某个客户端第一次建立同时包含恶意守卫与恶意出口的电路时，就将其判定为失守。

这个模型有意保持简单：忽略中间中继、家族、/16 网络、备用守卫和节点故障。为了保证速度，它没有使用每次抽样都重新计算数千个权重的 \texttt{random.choices}，而是在累积权重数组上进行二分搜索，从而无需每次都重新计算全部权重；具体运行时间取决于机器。

要复现这些数字，请先获取公开的历史共识（如执行 \texttt{make -C labs/zh-Hans data}）；该文件未直接存入 Git。模拟是确定性的：在模型、共识、比例和随机种子都相同的情况下，结果总是如下：
\begin{lstlisting}[language=bash,numbers=none]
$ PYTHONPATH=labs/zh-Hans/代码 python3 -m torlab.simulation labs/zh-Hans/数据/collector/consensus-microdesc-2026-09-23-15 0.05
攻击者：守卫带宽 5%、出口带宽 5%；2000 名客户端、180 天、每天 20 条电路
持久守卫             1d:  2.6%    7d:  4.4%   30d:  4.4%   90d:  9.4%  180d: 18.4%  中位数：None
每条电路重新选守卫   1d:  4.3%    7d: 27.9%   30d: 76.5%   90d: 98.9%  180d:100.0%  中位数：15
\end{lstlisting}

\begin{figure}[htbp]
\centering
% SPDX-License-Identifier: GPL-3.0-or-later
% Copyright (C) 2026 Alessandro Greco
%% EPUB 无障碍替代文本。
\altfig{实验 19 的两条累计曲线：教学模拟包含 2000 名客户端、180 天、每人每天 20 条电路。攻击者分别控制 5\% 的守卫和出口加权带宽。模型中的失守指至少一次使用了守卫和出口均由攻击者控制的电路；到第 180 天，每条电路换守卫的比例为 100\%，持久守卫为 18.4\%。各客户端按模拟的 48--60 天间隔分别轮换守卫，并非同步发生。这不是实际网络去匿名化观测结果。}
% 各客户端的轮换相互独立（模拟间隔 48--60 天），不能标记同步日期。
\begin{tikzpicture}[font=\sffamily\scriptsize,x=0.04cm,y=0.04cm]
  \draw[->] (0,0) -- (190,0) node[below left] {天};
  \draw[->] (0,0) -- (0,108) node[above] {模型失守客户端（\%）};
  \foreach \x in {0,30,60,90,120,150,180} \draw (\x,0) -- (\x,-2) node[below] {\x};
  \foreach \y in {0,25,50,75,100} \draw (0,\y) -- (-2,\y) node[left] {\y};
  \draw[thick,avviso] plot coordinates {(0,0.0) (5,21.1) (10,38.5) (15,51.1) (20,61.5) (25,69.7) (30,76.5) (35,81.7) (40,86.0) (45,89.2) (50,91.7) (55,93.3) (60,94.9) (65,96.2) (70,97.2) (75,97.9) (80,98.2) (85,98.7) (90,98.9) (95,99.1) (100,99.5) (110,99.7) (120,99.9) (165,100.0) (180,100.0)};
  \draw[thick,torviola] plot coordinates {(0,0.0) (1,2.6) (3,4.3) (5,4.4) (48,4.4) (52,5.8) (56,7.1) (60,8.7) (64,9.4) (100,9.5) (108,10.8) (114,12.4) (120,13.2) (130,13.4) (160,14.4) (165,15.6) (170,17.0) (175,18.2) (180,18.4)};
  \node[avviso,anchor=west] at (40,80) {每条电路使用新守卫};
  \node[torviola,anchor=west] at (66,26) {持久守卫};
% 各客户端的轮换相互独立（模拟间隔 48--60 天），不能标记同步日期。
  \node[torgrigio,anchor=west,align=left] at (66,43) {守卫轮换\\并不同步};
\end{tikzpicture}

\caption{实验~19 的简化模拟：使用 2026 年 9 月 23 日的共识快照，模拟 2000 名客户端在 180 天内每天建立 20 条电路；攻击者分别控制 5\% 的守卫和出口加权带宽。累计比例表示在模型中至少有一条电路同时使用恶意守卫与恶意出口的客户端。持久守卫轮换根据客户端独立发生（模拟间隔 48--60 天），不存在全体同步的三个轮换日期。这不是公共网络真实去匿名化事件的观测值。}
\label{fig:compromissione}
\end{figure}

\Cref{fig:compromissione} 显示了两条曲线。如果不使用持久守卫，十五天内就有一半客户端失守，三个月内几乎全部失守，因为每条新电路都提供一次新的独立机会。使用持久守卫时，风险主要在选择守卫的时刻确定。约 4\% 的客户端早早选择了恶意守卫并失守；其余客户端在轮换守卫前受到保护。在该模型中，各客户端约每 48--60 天分别轮换守卫、并非同步发生，新的抽样机会使更多客户端暴露，六个月后总失守比例达到约 18\%。如果攻击者控制的带宽比例降至 1\%，持久守卫情况下的六个月失守比例约为 4\%，而不使用持久守卫时约为 31\%。

这些基于 2026 年 9 月参考快照的结果，说明了守卫设计的理由。在这一简化模型中，守卫的主要作用是改变不同客户端之间的风险分布，而不是保证每条电路的风险都会下降。模型中，大多数用户一开始的风险接近零，但少数用户的风险很高。因此，守卫寿命是一种权衡：轮换越慢，使用诚实守卫的用户就能受到更长时间的保护，但选中恶意守卫的用户也会更长时间处于暴露状态。能够在网络中潜伏数月的攻击者，每次轮换都会获得新机会。

% ---------------------------------------------------------------------------
\section{如何安全地开展研究}\label{sec:ric-sicurezza}

\Cref{cap:25-legale} 介绍了 Tor Research Safety Board 的指南~\cite{trsb}。设计实验时，首先应当问：\emph{我的研究问题真的需要公共网络吗？}许多性能和攻击问题，都可以用 Shadow 或私有网络研究，并得到可重复的实验和可控的变量。当公共网络测量确实不可替代时，应遵循前面介绍的原则：只针对自己产生的流量操作，只收集最少的数据，在发布前聚合结果，并在开始前咨询安全委员会和 Tor Project。

近期历史提供了形成鲜明对比的案例。\cref{cap:27-casi-rete} 描述了一种通过接收中继可见、但外部观察者无法直接从加密 TLS 流量中读取的信号影响真实用户的实验或攻击。PrivCount 则表明，可以在不暴露个体用户的前提下细致测量网络。两者的区别不仅在于工具，更在于实验设计。

% ---------------------------------------------------------------------------
\section*{本章小结}
\begin{itemize}
  \item Tor 研究常发表于 PETS 及重要安全会议。Tor Project 提供数据、工具和 Research Safety Board。
  \item 熵衡量单次观察下的匿名性，匿名程度则把结果标准化到零与一之间。对长时间风险，还需要特定时间段内失守概率等指标。
  \item 分析模型、行为模拟、网络仿真或模拟，以及公共网络测量，分别用于回答不同问题，也具有不同风险。
  \item 基于参考共识的模拟显示，持久守卫通过将暴露风险集中到较少客户端身上，显著降低了随时间失守的客户端比例。
\end{itemize}

% ---------------------------------------------------------------------------
\section*{练习}

\begin{esercizio}[1]\label{es:29-entropia}
某攻击者观察一个包含四名用户的系统，赋予各用户的发送概率依次为 $1/2$、$1/4$、$1/8$ 和 $1/8$。计算熵、匿名程度，以及具有相同不确定性的等概率用户数量。
\end{esercizio}

\begin{esercizio}[1]\label{es:29-metodo}
为每个问题选择最合适的研究方法并说明理由：（a）如果电路使用四跳，延迟会如何变化？（b）某个国家有多少用户访问 Tor？（c）一种新的守卫选择算法能否延长首次失守时间？
\end{esercizio}

\begin{esercizio}[2]\label{es:29-durata}
运行实验~19 的 \texttt{simulate}，分别将 \texttt{guard\_lifetime} 从默认 60 天改为 30 天和 180 天，比较 180 天后客户端失守比例，并解释结果。
\end{esercizio}

\begin{esercizio}[2]\label{es:29-asimmetria}
重复实验~19：先假设攻击者控制 10\% 的出口带宽，却不控制任何守卫带宽；然后假设控制 10\% 的守卫带宽，却不控制出口带宽。各自能实现什么效果？为什么？
\end{esercizio}

\begin{esercizio}[3]\label{es:29-torps}
为模拟器增加目前忽略的三项功能：客户端对电路复用十分钟而不是不断新建电路；守卫与出口之间的 /16 约束；攻击者控制某大型网络所属的整个自治系统。哪些修改对结果影响最大？
\end{esercizio}

% ---------------------------------------------------------------------------
\section*{延伸阅读}
\begin{itemize}
  \item C.~Díaz、S.~Seys、J.~Claessens、B.~Preneel，\titolo{Towards Measuring Anonymity}~\cite{diaz2002measuring}。
  \item R.~Jansen、A.~Johnson，\titolo{Safely Measuring Tor}~\cite{jansen2016privcount}。
  \item Tor Project 研究门户：\url{https://research.torproject.org/}。
  \item \emph{Anonymity Bibliography}：\url{https://www.freehaven.net/anonbib/}。
\end{itemize}
